% ---------------------------------------------------------------- motivation & setup
\subsection{Motivation and protocol}
\label{sec:os-setup}

Measured against a model trained on the target building at matched sample size, the
cross-building penalty for activity recognition in this pair of apartments is $0.6$
macro-F1 points entering $\text{S}_2$ and $21.3$ entering $\text{S}_1$: in one direction
there is essentially no distribution shift for an adaptation method to recover, and in the
other there is a substantial one. That asymmetry, rather than any single method's score,
is what decided the outcome of our closed-set study, and the methods tracked it --- losing
ground wherever there was no headroom and gaining only where there was.

This section asks whether a second, simultaneous shift changes that. We withhold two
activities, \texttt{toileting} and \texttt{cleaning}, from source training and treat them
as target-private, so the model faces a change of building and a change of label space at
once. The retained classes are \texttt{meal\_prep}, \texttt{watching\_tv} and
\texttt{music}; \texttt{cleaning} is placed in the unknown set deliberately, since at $113$
target windows it is not learnable in this data and leaving it among the known classes
would confound open-set performance with a sample-size limit. The experiment is a
$2\times2$ grid, $\{\text{S}_1\!\to\!\text{S}_2, \text{S}_2\!\to\!\text{S}_1\}$ crossed
with $\{$closed-set, open-set$\}$, and because withholding classes changes the source
model, every cell is given its own penalty measurement before any method is run.

Throughout we hold to the Tier-1 definition of test-time adaptation used in the
closed-set study: the source model is trained once and then frozen, with no further
multi-epoch training on target data; adaptation is online, in a single pass over the
unlabelled target stream, in recording order, with each batch predicted before the model
is updated on it; source data is unavailable during adaptation; and no target labels are
used anywhere, including for selecting hyperparameters and for calibrating the
unknown-rejection threshold, which is fixed from source validation before the target
stream is read. Any quantity a method relies on --- scalers, prototypes, thresholds ---
must be estimable without target labels. Where we report a figure that violates this, as
with the labelled ceiling of Section~\ref{sec:os-depth} and the matched-rejection-rate
comparison of Section~\ref{sec:os-maha}, it is labelled a diagnostic and is used to bound
what a method could achieve rather than to claim what one did.

The constraint matters because the reference implementations in this area do not observe
it. Section~\ref{sec:os-lead} documents four places where the released code for LEAD
leaves the test-time setting, including fifty offline epochs over the target set and an
evaluation loader that wraps the same data the model adapted on. Results obtained that way
are not comparable with results obtained under the constraints above, which is why every
number in this section is reported against our own Tier-1 source-only baseline rather than
against published figures.
